\documentclass[10pt,a4paper]{amsart}
\usepackage[margin=19mm]{geometry}
\usepackage{amsmath,amssymb,mathtools}
\usepackage[scr=rsfs,cal=euler]{mathalfa}
\usepackage{xfrac}
\usepackage[unicode,hidelinks,bookmarks=false]{hyperref}
\newcommand{\ie}{i.e.\ }
\newcommand{\cf}{cf.\ }
\newcommand{\resp}{respectively\ }
\newcommand{\Subsection}{\subsection}
\providecommand{\nobreakdash}{\nobreak\hbox{-}}
\setlength{\emergencystretch}{2em}
\multlinegap0pt
\usepackage{amssymb}
\usepackage[shortlabels]{enumitem}
\newtheorem{lemma}{Lemma}
\renewcommand{\Re}{\operatorname{Re}}
\newcommand\matn{{\mathcal M}_n}
\newcommand\ta{{\tilde{A}}}
\title{From Clou\^atre-Ostermann-Ransford to Okubo-Ando}
\author{Michael Hartz, John E.\ M\textsuperscript{c}Carthy}
\date{Source: arXiv:2606.02922v1 (CC BY 4.0)}
\begin{document}
\maketitle

\noindent\emph{Source and licence.} From Clou\^atre-Ostermann-Ransford to Okubo-Ando. Version arXiv:2606.02922v1 (CC BY 4.0), distributed by arXiv under the Creative Commons Attribution 4.0 International licence, http://creativecommons.org/licenses/by/4.0/, as recorded in the arXiv licence metadata for this exact version and shown on the abstract page https://arxiv.org/abs/2606.02922v1. Full licence notice: \texttt{licenses/arxiv-2606.02922-CC-BY-4.0.txt}.\par\smallskip

\noindent\textit{Editorial scope.} Counted unit: the article's lemma lem:extremal (source0.tex lines 294--307). The article's complete original proof of it is delivered with it (source0.tex lines 309--355, one \texttt{proof} environment of 47 lines). No mathematics is written, repaired or re-derived: the statement and the proof are the article's own lines, in the article's own notation, and no retained line is substituted or re-flowed.  No further source-local context is needed: every cross-reference of every retained line resolves inside the two delivered spans.  Nothing was omitted from any retained span, so the package carries no omission record.  No retained line contains a citation, so the wrapper carries no bibliography and no bibliography entry.  No retained line contains a figure environment, an image-inclusion command or drawing code, so no image is delivered and the package declares no asset dependency.  Editorial formatting only, in addition: an AMS \texttt{amsart} preamble that restates the article's own declarations and macros used by the retained lines, namely the environments \texttt{lemma} on separate counters, together with the standard packages those lines need.  The retained environments print in this document as Lemma 1 in order of appearance, whereas the article prints the same items with the numbering of its own full document.  The complete native source of the article as distributed in the arXiv e-print for this exact version is bundled unchanged as \texttt{sources/201941/source0.tex}.
\begin{lemma}
  \label{lem:extremal}
  Let $(R_k)$ be a sequence in $\matn(B(\mathcal{H}))$. Let $C \in (0,\infty) \setminus \{1\}$.
  Assume that whenever $A \in \matn(\mathbb{C})$ satisfies $\|A\| \|R_k\| < 1$ then $ \|m_\ta(R_k)\| \le C$.
%  \begin{equation*}
%    \|m_A(R_k)\| \le C \text{ for all } k \in \mathbb{N} \text{ and } A \in \matn(\mathbb{C}) \text{ with } \|A\| \|R_k\| < 1.
%  \end{equation*}

  Let $(x_k)$ be a sequence of unit vectors in $\mathcal{H}^n$ with $\lim_{k \to \infty} \|R_k x_k\| = C$.
  Then for every bounded sequence $(A_k)$ in $\matn(\mathbb{C})$,
  \begin{equation*}
    \lim_{k \to \infty} \langle R_k x_k, (A_k \otimes I_{\mathcal{H}}) x_k \rangle  =0.
  \end{equation*}
\end{lemma}

\begin{proof}
  The assumption implies that $\|R_k\| \le C$ for all $k$ by taking $A=0$.
  Let $A \in \matn(\mathbb{C})$ with $C \|A\| < 1$ and let $y \in \mathcal{H}^n$ be a unit vector.
 % To shorten notation, we identify $A$ with $A \otimes I_{\mathcal{H}}$.
  Define $z_k = (I - \ta^* \ta)^{-1/2} (I + \ta^* R_k) y$.
  Then by assumption,
  \begin{align*}
    \|(I - \ta^* \ta)^{-1/2} (R_k + \ta) y\|^2  &= \|m_\ta(R_k) z_k\|^2 \\ &\le C^2 \|z_k\|^2  \\
                                                        &= C^2 \|(I - \ta^* \ta)^{-1/2} (I + \ta^* R_k) y\|^2.
  \end{align*}

  We now replace $A$ with $t A$ for a small complex number $t$ and expand to first order.
  We have $(I - |t|^2 \ta \ta^*)^{-1/2} = I + \mathcal{O}(|t|^2)$ and $(I - |t|^2 \ta^* \ta)^{-1/2} = I + \mathcal{O}(|t|^2)$.
  Therefore, the last inequality implies that
  \begin{equation*}
    \|R_k y\|^2 + 2 \Re \overline{t} \langle R_k y, \ta y \rangle \le C^2 ( \|y\|^2 + 2 \Re t \langle y, \ta^* R_k y \rangle )
    + \mathcal{O}(|t|^2),
  \end{equation*}
  where the implied constant in the $\mathcal{O}(|t|^2)$ term is independent of the unit vector $y$ and the number $k$
  because $\|R_k\| \le C$.
  Rearranging yields
  \begin{equation*}
 2   (1 - C^2)  \Re \overline{t} \langle R_k y , \ta y \rangle \le C^2 \|y\|^2 - \|R_k y\|^2 + \mathcal{O}(|t|^2)
  \end{equation*}
  for all complex numbers $t$ in a neighborhood of zero and all unit vectors $y$.
  By varying the argument of $t$, we conclude that
  \begin{equation*}
   2 |1 - C^2|  |t| | \langle R_k y , \ta y \rangle| \le C^2 \|y\|^2 - \|R_k y\|^2 + \mathcal{O}(|t|^2).
  \end{equation*}

  We apply the last inequality to $y = x_k$, the sequence in the statement of the lemma. Since the estimate
  above is independent of the unit vector $y$ and $k$, it follows that
  \begin{equation*}
 2   |1 - C^2|  |t| \limsup_{k \to \infty} | \langle R_k x_k , \ta x_k \rangle| \le \mathcal{O}(|t|^2)
  \end{equation*}
  for all $t$ in a neighborhood of zero. Since $C \neq 1$, we conclude that
  \[
    \lim_{k \to \infty} \langle R_k x_k, \ta x_k \rangle = 0.
  \]
  We have proved this identity for all $A \in \matn(\mathbb{C})$ with $ C \|A\| < 1$,
  but then by homogeneity, it holds for all $A \in \matn(\mathbb{C})$. This shows that
  the linear functionals
  \begin{equation*}
    \matn(\mathbb{C}) \to \mathbb{C}, \quad A \mapsto \langle \ta x_k, R_k x_k \rangle,
  \end{equation*}
  converge to zero pointwise, and hence in norm since $\matn(\mathbb{C})$ is finite dimensional. This completes the proof.
\end{proof}

\end{document}
